\section*{الفصل \ref{ch:g12:internal-heat} --- الحرارة الداخلية للأرض}

\begin{solution}{exo:g12:internal-heat:1}
(أ) \omterm{def:g12:internal-heat:heat-flow}{التدفق الحراري}؛ (ب) \omterm{def:g12:internal-heat:heat-production}{إنتاج الحرارة الإشعاعية}؛ (ج) \omterm{def:g12:internal-heat:secular-cooling}{التبرد القرني}؛
(د) \omterm{def:g12:internal-heat:convection}{الحمل الحراري}؛ (ه) \omterm{def:g12:internal-heat:geothermal}{الطاقة الحرارية الأرضية}.
\end{solution}

\begin{solution}{exo:g12:internal-heat:2}
(1) \qty{0.0916}{W/m^2}. (2) \qty{4.67e13}{W}. (3) \qty{0.029}{K/m}: فالفرق
بدرجة سلسيوس واحدة فرق بكلفن واحد.
\end{solution}

\begin{solution}{exo:g12:internal-heat:3}
الثوريوم-232، واليورانيوم-238، والبوتاسيوم-40، واليورانيوم-235. ولا يرصد كاملاند
إلا النيوترينوات الأرضية للثوريوم-232 واليورانيوم-238.
\end{solution}

\begin{solution}{exo:g12:internal-heat:4}
$490/\sqrt{1} = \qty{490}{mW/m^2}$؛ $490/2 = \qty{245}{mW/m^2}$؛
$490/5 = \qty{98}{mW/m^2}$؛ $490/7 = \qty{70}{mW/m^2}$، وهذا قريب من متوسط
القارات.
\end{solution}

\begin{solution}{exo:g12:internal-heat:5}
قرب السطح، تتبع درجة حرارة الأرض الأيام
والفصول؛ ولا تثبت عند متوسط درجة حرارة السنة إلا تحت بضعة أمتار،
ولا تُظهر الحرارة الآتية من الأسفل إلا في عمق أكبر من ذلك.
والحفر يحرك الصخر والماء في البئر ويغيّر
درجة حرارتهما، فننتظر حتى يعود البئر إلى درجة حرارة
الصخر.
\end{solution}

\begin{solution}{exo:g12:internal-heat:6}
(1) $(265 - 10)/9.1 \approx \qty{28.0}{\celsius}$ لكل كيلومتر، أي
\qty{0.0280}{K/m}. (2) $q = 2.72 \times 0.0280 \approx \qty{0.0762}{W/m^2}$،
أي نحو \qty{76}{mW/m^2}. (3) فوق متوسط القارات بقليل،
\qty{70.9}{mW/m^2}، بنحو \qty{7}{\percent}: وهي قيمة قارية نموذجية.
\end{solution}

\begin{solution}{exo:g12:internal-heat:7}
القارات: $0.0709 \times \num{2.073e14} \approx \qty{1.47e13}{W} =
\qty{14.7}{TW}$. والمحيطات: $0.1054 \times \num{3.028e14} \approx
\qty{3.19e13}{W} = \qty{31.9}{TW}$. والمجموع نحو \qty{46.6}{TW}، يغادر منه
$31.9/46.6 \approx \qty{68}{\percent}$، أي الثلثان، عبر قاع
المحيط.
\end{solution}

\begin{solution}{exo:g12:internal-heat:8}
(1) $\num{1.88e-6} \times 8000 \approx \qty{0.0150}{W/m^2} = \qty{15.0}{mW/m^2}$.
(2) $15.0/70.9 \approx \qty{21}{\percent}$، أي خُمس \omterm{def:g12:internal-heat:heat-flow}{التدفق الحراري}.
\end{solution}

\begin{solution}{exo:g12:internal-heat:9}
$\sqrt{t} = 490/70.9 \approx 6.91$، إذن $t \approx 47.8$ مليون سنة: فقاع
المحيط الأقدم من نحو 48 مليون سنة يفقد من الحرارة أقل مما تفقده القارة
المتوسطة.
\end{solution}

\begin{solution}{exo:g12:internal-heat:10}
اليورانيوم-235: $2^{1.4/0.704} = 2^{1.99} \approx 4$ مرات أكثر. والبوتاسيوم-40:
$2^{2.5/1.248} = 2^{2.00} \approx 4$ مرات أكثر. فكلاهما رجوع بعمرَي نصف.
\end{solution}

\begin{solution}{exo:g12:internal-heat:11}
(1) $10^2/10^{-6} = \qty{1e8}{s}$، أي نحو 3 سنوات. (2) $(10^5)^2/10^{-6} =
\qty{1e16}{s}$، أي نحو 320 مليون سنة. (3) $(\num{2.891e6})^2/10^{-6}
\approx \qty{8.4e18}{s}$، أي نحو 265 مليار سنة. فالزمن يزداد بمربع
المسافة: سمك أكبر بعشر مرات، وزمن أطول بمئة مرة.
\end{solution}

\begin{solution}{exo:g12:internal-heat:12}
اضرب كلًّا منها في $2^{3/T}$: الثوريوم-232، $10.4 \times 2^{3/14.0} \approx
\qty{12.1}{TW}$؛ واليورانيوم-238، $9.6 \times 2^{3/4.468} \approx \qty{15.3}{TW}$؛
والبوتاسيوم-40، $4.0 \times 2^{3/1.248} \approx \qty{21.2}{TW}$؛ واليورانيوم-235،
$0.42 \times 2^{3/0.704} \approx \qty{8.1}{TW}$. والمجموع نحو \qty{56.7}{TW}،
أي أكثر من ضعف \qty{24.4}{TW} اليوم. ويعطي النظيران القصيرا العمر
$(21.2 + 8.1)/56.7 \approx \qty{52}{\percent}$ منه، مقابل
$4.42/24.4 \approx \qty{18}{\percent}$ اليوم.
\end{solution}

\begin{solution}{exo:g12:internal-heat:13}
القدرة هي الكتلة في الثانية مضروبة في الحرارة النوعية مضروبة في
الارتفاع في درجة الحرارة: $\num{3e8} = 1640 \times 4180 \times \Delta T$، إذن $\Delta T =
\num{3e8}/(1640 \times 4180) \approx \qty{44}{K}$. والقارة المتوسطة تفقد
\qty{0.0709}{W} لكل متر مربع، فإن \qty{300}{MW} هي الفقد الطبيعي للحرارة
من $\num{3e8}/0.0709 \approx \qty{4.2e9}{m^2}$، أي نحو \qty{4200}{km^2}:
فالمحطة تجمع حرارة تُخرجها الأرض طبيعيًا على آلاف
الكيلومترات المربعة.
\end{solution}

\begin{solution}{exo:g12:internal-heat:14}
(1) 100 مليون سنة أقصر من 260 مليار سنة بمقدار $\num{2.6e11}/\num{1e8} = 2600$
مرة؛ و200 مليون سنة أقصر بمقدار 1300 مرة.
(2) الحمل ينقل الحرارة مع الصخر المتحرك: فزمن عبور
الوشاح هو سمكه مقسومًا على السرعة، $\qty{2891}{km}$ ببضعة
سنتيمترات في السنة، أي نحو مئة مليون سنة. أما التوصيل فلا ناقل له
كهذا: فزمنه يزداد بمربع المسافة، وعلى آلاف
الكيلومترات يصير هائلًا.
\end{solution}

\begin{solution}{exo:g12:internal-heat:15}
(1) $10/22.3 \approx \qty{45}{\percent}$؛ والباقي تبرد
الوشاح نفسه، الذي يفقد حرارته المختزنة فيه. (2) يقع الحد على نصف قطر $6371 - 2891 =
\qty{3480}{km}$؛ ومساحته $4\pi (\num{3.48e6})^2 \approx
\qty{1.52e14}{m^2}$، و$\num{1e13}/\num{1.52e14} \approx
\qty{0.066}{W/m^2}$، أي نحو \qty{66}{mW/m^2}، وهذا نحو \qty{72}{\percent} من
المتوسط عند السطح. (3) تعبر الحرارة الحد بالتوصيل إلى
أسفل طبقة في الوشاح؛ فيصير هذا الصخر، إذ يسخن، أخف من
الوشاح الذي فوقه ويستطيع أن يصعد حاملًا الحرارة إلى الأعلى: وهكذا يُعتقد أن الأعمدة
الصاعدة تبدأ، في الأعماق.
\end{solution}

\begin{solution}{pb:g12:internal-heat:1}
\textbf{1.} $4\pi (\num{6.371e6})^2 \approx \qty{5.101e14}{m^2}$،
و$2.073 + 3.028 = 5.101$: نعم.

\textbf{2.} $3.028/5.101 \approx \qty{59}{\percent}$، أي نحو ثلاثة أخماس.

\textbf{3.} قاع المحيط حديث: فهو يُصنع ساخنًا عند الظهرات
ويبرد كلما ابتعد، \omterm{def:g11:lithosphere-asthenosphere:lithosphere}{والغلاف الصخري} الحديث يفقد حرارة كثيرة.

\textbf{4.} $490/\sqrt{4} = \qty{245}{mW/m^2}$.

\textbf{5.} $2.72 \times 0.026 \approx \qty{0.0707}{W/m^2} =
\qty{70.7}{mW/m^2}$: دون المتوسط بقليل، ويكاد يساويه.

\textbf{6.} $0.0709 \times \num{2.073e14} \approx \qty{14.7}{TW}$.

\textbf{7.} $0.1054 \times \num{3.028e14} \approx \qty{31.9}{TW}$.

\textbf{8.} $14.7 + 31.9 \approx \qty{46.6}{TW}$، وهي \qty{47}{TW} في
التقدير المنشور (والمتوسط العالمي في الفصل يعطي \qty{46.7}{TW}، وهي
القيمة المستعملة فيما يلي).

\textbf{9.} $\num{4.66e13}/\num{5.101e14} \approx \qty{0.0914}{W/m^2}$،
أي نحو \qty{91}{mW/m^2}.

\textbf{10.} $0.0243 \times \num{2.073e14} \approx \qty{5.0}{TW}$، أي ثلث
\qty{14.7}{TW} المفقودة عبر القارات: فالنشاط الإشعاعي في
\omterm{def:g12:continental-crust:continental-crust}{القشرة القارية} جزء كبير من تدفقها الحراري.

\textbf{11.} لكل كيلوغرام من اليورانيوم تحوي الأرض \qty{3.9}{kg} من
الثوريوم، تحرر $3.9 \times 26.3 \approx \qty{102.6}{\micro W}$،
مقابل \qty{94.4}{\micro W} لليورانيوم-238. فاليورانيوم-238 يصنع إذن
$94.4/(94.4 + 102.6) \approx 0.479$ من \qty{20.0}{TW}: أي نحو
\qty{9.6}{TW}. واليورانيوم-235: $9.59 \times 4.10/94.4 \approx \qty{0.42}{TW}$.

\textbf{12.} $20.0 + 0.42 + 4 \approx \qty{24.4}{TW}$.

\textbf{13.} الغرانيت وسائر صخور \omterm{def:g12:continental-crust:continental-crust}{القشرة القارية}، وهي أغنى بنحو
عشرين مرة من الصخور فوق المافية للوشاح.

\textbf{14.} أكبر: فكل نظير كان أوفر حينئذ، ضعف الكمية عن
كل عمر نصف إلى الوراء؛ وقبل 3 مليارات سنة كان المجموع أكثر من ضعف
مجموع اليوم.

\textbf{15.} $46.7 - 24.4 = \qty{22.3}{TW}$.

\textbf{16.} من الحرارة المختزنة منذ تكوّن الأرض: أي \omterm{def:g12:internal-heat:secular-cooling}{التبرد
القرني} للوشاح واللب.

\textbf{17.} $22.3/46.7 \approx \qty{48}{\percent}$.

\textbf{18.} من $11/46.7 \approx \qty{24}{\percent}$ إلى $38/46.7 \approx
\qty{81}{\percent}$: والنيوترينوات الأرضية تضيّق هذا المدى كثيرًا.

\textbf{19.} \omterm{def:g12:internal-heat:convection}{الحمل الحراري في الوشاح}، الذي ينقل معظم الحرارة؛ ولا تُعبر بالتوصيل
إلا العشرات الأخيرة من الكيلومترات.

\textbf{20.} $24.4/46.7 \approx 0.52$: فالنشاط الإشعاعي يوفر
\textbf{نحو نصف} ما تفقده الأرض من حرارة اليوم، ويوفر التبرد البطيء
للكوكب النصف الآخر.
\end{solution}
